% ch5.tex -- Module 5, Chapter 5: Energy in gravitational fields
% Syllabus: M5.IQ3.05 (primary); links to M5.IQ3.03, M5.IQ3.04 (Chapter 4) and M5.IQ2.04.
% ch5_writer_numbers_output.txt) and ch5_orbit_sim.py (Key experiment data).
% Data-sheet constants: G = 6.67e-11, M_E = 6.0e24 kg, r_E = 6.371e6 m, g = 9.8.
\chapter[Energy in Gravitational Fields]{Energy in Gravitational Fields: Orbits, Escape and Changing Orbits}\label{ch5:chap}
\covers{M5.IQ3.05}

\inquiry{How does the force of gravity determine the motion of planets and satellites?}

\begin{goals}
  \item explain why gravitational potential energy is negative, and use
        $U = -\dfrac{GMm}{r}$
  \item derive and use $K = \dfrac{GMm}{2r}$ and $U + K = -\dfrac{GMm}{2r}$ for a
        satellite in a circular orbit
  \item explain, using $K$ and $U$, why a planet in an elliptical orbit is
        fastest at perihelion (link to Kepler's laws)
  \item derive and use the escape velocity $v_\mathrm{esc} = \sqrt{\dfrac{2GM}{r}}$
  \item calculate the energy needed to change orbit, and explain orbital decay
\end{goals}

\begin{syllabus}
  \item M5.IQ3.05: derive and apply gravitational force and potential energy in
        radial fields: escape velocity, $U = -GMm/r$, $U + K = -GMm/2r$, energy
        changes between orbits, Kepler's laws (with \cref{ch4:chap})
\end{syllabus}

%% ==================================================================
\section{Gravitational potential energy in a radial field}\label{ch5:sec:gpe}
\covers{M5.IQ3.05}

\begin{bigidea}
Far from the ground, gravity gets weaker, so $mg\Delta h$ stops working. We
choose the potential energy to be zero very far away. Then the potential energy
everywhere else is negative: $U = -\dfrac{GMm}{r}$.
\end{bigidea}

\begin{keywords}
  \kw{Gravitational potential energy ($U$)}{energy stored because of where an
    object is in a gravitational field. Unit: J.}
  \kw{Radial field}{a gravitational field (the region where a planet's gravity
    pulls) whose lines all point straight at the planet's centre, like the
    spokes of a wheel.}
  \kw{Infinity}{a place so far away that the planet's gravity there is too
    small to matter. (Further out, gravity gets closer and closer to zero.)}
\end{keywords}

\begin{refresher}{Work and energy from Year 11}
Work: $W = F_\parallel s = Fs\cos\theta$. Kinetic energy: $K = \dfrac{1}{2}mv^2$
(never negative). Near the ground: $\Delta U = mg\Delta h$. All three are on the
sheet. \term{Conservation of energy}: energy is never made or destroyed; it only
changes form.
\end{refresher}

\subsection{Work done by gravity, and why $U$ is negative}

Drop a ball. Gravity and the motion point the same way. So gravity does
\emph{positive} work, and the ball speeds up. Throw the ball up. Now gravity and
the motion point opposite ways. So gravity does \emph{negative} work, and the
ball slows down.

When gravity does positive work, the ball gains kinetic energy. That energy
comes out of its potential energy, so $U$ falls by the same amount. In symbols:
$W_\text{by gravity} = -\Delta U$.

To move an object away at a steady speed, an engine must push against gravity.
\term{Work against gravity} is the energy the engine gives the object to do
this. All of it is stored as potential energy: $W_\text{against} = +\Delta U$.

$\Delta U = mg\Delta h$ assumes that $g$ (the gravitational field strength,
\qty{9.8}{m.s^{-2}} at the surface) is the same at every height. But $g$
gets smaller as you go up: at \qty{1500}{km} it is only about
\qty{6.5}{m.s^{-2}}. So for large heights we need a better formula.

Potential energy needs a zero level. In space, the best zero is at infinity,
where gravity has died away to zero. So we choose $U = 0$ at $r = \infty$. Now follow the steps.
\begin{enumerate}
  \item Moving away from the planet needs work against gravity, so $U$
        \emph{increases} as $r$ increases.
  \item $U$ keeps increasing all the way out, and it reaches 0 only at infinity.
  \item So at every real distance, $U$ is less than zero: it is negative.
\end{enumerate}
Think of a deep well. Call the top of the well ``zero height''. Then every point
inside the well has a negative height. You need energy to climb out. A satellite
sits in Earth's ``gravity well'' in the same way.

The minus sign comes from our choice of zero. Only \emph{changes} in $U$ can be
measured.

\subsection{Where $U = -GMm/r$ comes from}

\begin{toolkit}{Work from the area under a graph}
$W = Fs$ works only when the force stays the same. When the force changes, the
work is the \emph{area} under the force--distance graph. Cut the area into thin
strips. Each strip is nearly a rectangle (height = force, width = distance), so
its area is force $\times$ distance. Then add up the strips.
\end{toolkit}

\begin{toolkit}{Subtracting fractions like $\frac{1}{a} - \frac{1}{b}$}
To subtract two fractions, give them the same bottom number (the
\emph{denominator}):
\[ \frac{1}{a} - \frac{1}{b} = \frac{b}{ab} - \frac{a}{ab} = \frac{b - a}{ab} \]
Example: $\tfrac{1}{2} - \tfrac{1}{4} = \tfrac{4 - 2}{8} = \tfrac{1}{4}$. Read
backwards, it splits one fraction into two:
$\tfrac{b - a}{ab} = \tfrac{1}{a} - \tfrac{1}{b}$. Warning:
$\tfrac{1}{a} - \tfrac{1}{b}$ is \emph{not} $\tfrac{1}{a - b}$.
\end{toolkit}

% final audit: strip edges are r_A, r_B (were italic r_a, r_b), so they cannot be confused with
% the apogee distance r_a (upright a) used from Section 5.2 on and in Chapter 4.
Gravity on a mass $m$ is $F = \dfrac{GMm}{r^2}$ (on the sheet). Take one thin
strip from $r_\mathrm{A}$ out to $r_\mathrm{B}$ (\cref{ch5:fig:FU}a). The force falls from
$\dfrac{GMm}{r_\mathrm{A}^2}$ to $\dfrac{GMm}{r_\mathrm{B}^2}$. The strip is very thin, so $r_\mathrm{A}$
and $r_\mathrm{B}$ are almost equal, and the force hardly changes across it. So any value
between the two ends is a good height for the strip.

\begin{figure}
  \centering
  \begin{tikzpicture}
    \begin{axis}[hsc, name=Fplot, width=0.50\linewidth, height=0.33\linewidth,
        xmin=1, xmax=6, ymin=0, ymax=11, xtick={1,2,...,6},
        xlabel={$r$ (in units of $r_\mathrm{E}$)},
        ylabel={$F$ on \qty{1}{kg} (\si{N})},
        title={(a) Force}, title style={font=\small\bfseries, text=hscink}]
      \addplot[draw=none, fill=hscseriesA!22, domain=1.5:3, samples=40] {9.86/x^2} \closedcycle;
      \addplot[hscseriesA, domain=1:6, samples=80, line width=1.2pt] {9.86/x^2};
      \node[font=\footnotesize, align=left, anchor=west, text=hscink] at (axis cs:2.7,6.2)
        {shaded area $= \Delta U$\\$=$ work done against\\gravity, $r_1$ to $r_2$};
      \draw[guide] (axis cs:1.5,0) -- (axis cs:1.5,4.38);
      \draw[guide] (axis cs:3,0) -- (axis cs:3,1.10);
      \node[font=\footnotesize, anchor=south west, inner sep=1pt] at (axis cs:1.53,0.15) {$r_1$};
      \node[font=\footnotesize, anchor=south east, inner sep=1pt] at (axis cs:2.97,0.15) {$r_2$};
    \end{axis}
    \begin{axis}[hsc, at={(Fplot.east)}, anchor=west, xshift=1.3cm,
        width=0.50\linewidth, height=0.33\linewidth,
        xmin=1, xmax=8, ymin=-70, ymax=10, xtick={1,2,...,8},
        xlabel={$r$ (in units of $r_\mathrm{E}$)},
        ylabel={$U$ of \qty{1}{kg} (\si{MJ})},
        title={(b) Potential energy}, title style={font=\small\bfseries, text=hscink}]
      \draw[hscink!70, line width=0.8pt] (axis cs:1,0) -- (axis cs:8,0);
      \addplot[hscseriesB, domain=1:8, samples=80, line width=1.2pt, solid] {-62.8/x};
      \node[font=\footnotesize, anchor=south, text=hscink] at (axis cs:4.6,0.5) {$U = 0$ at infinity};
      \node[font=\footnotesize, anchor=west, text=hscink] at (axis cs:1.15,-62.8) {surface: \qty{-62.8}{MJ}};
      \node[font=\footnotesize, align=left, anchor=west, text=hscink] at (axis cs:3.2,-36)
        {$U$ rises towards 0\\but never reaches it};
    \end{axis}
  \end{tikzpicture}
  \caption{(a) The area under the $F$--$r$ graph is the work done against
    gravity. (b) $U = -\dfrac{GMm}{r}$ for \qty{1}{kg} near Earth: negative
    everywhere, rising towards zero.}\label{ch5:fig:FU}
\end{figure}

We choose $\dfrac{GMm}{r_\mathrm{A} r_\mathrm{B}}$. It lies between the two ends, because
$r_\mathrm{A} r_\mathrm{B}$ lies between $r_\mathrm{A}^2$ and $r_\mathrm{B}^2$. It also makes the next step neat.
\begin{align*}
  W_\text{strip} &= \frac{GMm}{r_\mathrm{A} r_\mathrm{B}} \times (r_\mathrm{B} - r_\mathrm{A})
     && \text{height $\times$ width}\\
  &= GMm\left(\frac{1}{r_\mathrm{A}} - \frac{1}{r_\mathrm{B}}\right)
     && \text{split the fraction (Maths toolkit above)}
\end{align*}
Now cut the whole area from $r_1$ to $r_2$ into many thin strips. Call the edges
$r_1, r_\mathrm{A}, r_\mathrm{B}, \ldots, r_2$. Add the work for every strip:
\[ GMm\left(\frac{1}{r_1} - \frac{1}{r_\mathrm{A}}\right) + GMm\left(\frac{1}{r_\mathrm{A}} - \frac{1}{r_\mathrm{B}}\right) + \cdots \]
The $-\dfrac{1}{r_\mathrm{A}}$ and $+\dfrac{1}{r_\mathrm{A}}$ cancel, and so do all the other
inside terms. Only the two ends, $\dfrac{1}{r_1}$ and $\dfrac{1}{r_2}$, are left.
\begin{align}
  W_\text{against} = \Delta U &= GMm\left(\frac{1}{r_1} - \frac{1}{r_2}\right)
     && \text{work from $r_1$ out to $r_2$}\label{ch5:eq:dU}\\
  U(\infty) - U(r) &= GMm\left(\frac{1}{r} - \frac{1}{\infty}\right)
     && \text{start at $r_1 = r$, go all the way to $r_2 = \infty$}\notag\\
  0 - U(r) &= GMm\left(\frac{1}{r} - 0\right)
     && U(\infty) = 0 \text{ and } \tfrac{1}{\infty} = 0\notag\\
  U(r) &= -\frac{GMm}{r}
     && \text{multiply both sides by $-1$}\notag
\end{align}
(With calculus: $\int_r^\infty \frac{GMm}{x^2}\,\mathrm{d}x = \frac{GMm}{r}$. You
do not need calculus for the HSC.)

\textbf{Link to Year 11.} Put $r_1 = r_\mathrm{E}$ and $r_2 = r_\mathrm{E} + h$
into \cref{ch5:eq:dU}:
\[ \Delta U = GMm\left(\frac{1}{r_\mathrm{E}} - \frac{1}{r_\mathrm{E} + h}\right)
   = GMm\,\frac{(r_\mathrm{E} + h) - r_\mathrm{E}}{r_\mathrm{E}(r_\mathrm{E} + h)}
   = GMm\,\frac{h}{r_\mathrm{E}(r_\mathrm{E} + h)} \]
For a small $h$, $r_\mathrm{E} + h \approx r_\mathrm{E}$, so
$\Delta U \approx \dfrac{GM}{r_\mathrm{E}^2}\,mh$. Since
$g = \dfrac{GM}{r_\mathrm{E}^2}$, this is $mgh$.

But $mg\Delta h$ with $g = \qty{9.8}{m.s^{-2}}$ is 4\,\% too big at
\qty{300}{km} and 23\,\% too big at \qty{1500}{km}. At geostationary height
(about \qty{36000}{km} up) it gives 6.6 times the true value.

\begin{formulacard}{Gravitational potential energy}
  \cardformula{U = -\frac{GMm}{r}}
  \begin{symbols}
    \sym{$U$}{gravitational potential energy (zero at infinity)}{\si{J}}
    \sym{$G$}{universal gravitational constant, \num{6.67e-11}}{\si{N.m^2.kg^{-2}}}
    \sym{$M$}{mass of the planet or star}{\si{kg}}
    \sym{$m$}{mass of the object}{\si{kg}}
    \sym{$r$}{distance from the \emph{centre} of the planet}{\si{m}}
  \end{symbols}
  \inwords{the closer an object is to the planet, the lower (more negative) its $U$.}
  \whentouse{at any distance from a planet, moon or star: surface, orbit or far away.}
  \whennottouse{inside the planet ($r$ smaller than the planet's radius).}
  \rearranged{$r = -\dfrac{GMm}{U}$. Between two distances (combine two values of $U$):
    $\Delta U = GMm\left(\dfrac{1}{r_1} - \dfrac{1}{r_2}\right)$, $W_\text{by gravity} = -\Delta U$.}
  \tensecond{\qty{1.0}{kg} at $r = \qty{1.0e7}{m}$ from Earth's centre:
    $U = -\dfrac{(\num{6.67e-11})(\num{6.0e24})(1.0)}{\num{1.0e7}} = \qty{-4.0e7}{J}$.}
  \trap{$r = r_\mathrm{E} + \text{altitude}$. Keep the minus sign. Never use
    $\dfrac{GMm}{r_2 - r_1}$ for $\Delta U$: subtract the two values of $U$.}
\end{formulacard}

\begin{toolkit}{Ratio questions: what happens when $r$ doubles?}
If $y = \dfrac{k}{r}$ and $r$ doubles, $y$ becomes $\dfrac{k}{2r} = \tfrac{1}{2}y$.

If $y = \sqrt{\dfrac{k}{r}}$ and $r$ doubles, $y$ becomes
$\tfrac{1}{\sqrt{2}} = 0.71$ times as big.

If $y$ is negative, halving it moves it closer to zero, so $y$
\emph{increases}. Example: \qty{-8}{J} becomes \qty{-4}{J}.
\end{toolkit}

\begin{workedexample}{Potential energy of a weather satellite}\label{ch5:we:weather}
A \qty{720}{kg} weather satellite orbits \qty{850}{km} above Earth's surface.
Calculate its gravitational potential energy.
\given{$m = \qty{720}{kg}$, \ altitude $= \qty{850}{km} = 850 \times \qty{1000}{m} = \qty{8.50e5}{m}$, \
  $M_\mathrm{E} = \qty{6.0e24}{kg}$, \ $r_\mathrm{E} = \qty{6.371e6}{m}$}
\find{$U$ (first find $r$ from Earth's centre)}
\formula{$r = r_\mathrm{E} + \text{altitude}$ \quad then \quad $U = -\dfrac{GMm}{r}$}
\substitute{\eqlines{r &= \qty{6.371e6}{m} + \qty{0.850e6}{m} = \qty{7.221e6}{m}\\
  U &= -\frac{(\qty{6.67e-11}{N.m^2.kg^{-2}})(\qty{6.0e24}{kg})(\qty{720}{kg})}{\qty{7.221e6}{m}}\\
     &\quad\text{(N m$^2$ kg$^{-2}$ $\times$ kg $\times$ kg = N m$^2$; and N m = J, so N m$^2$ = J m)}\\
     &= -\frac{\qty{2.881e17}{J.m}}{\qty{7.221e6}{m}} = \qty{-3.990e10}{J}
        \quad\text{(J m $\div$ m = J)}}}
\answer{\final{$U = \qty{-3.99e10}{J}$} ($M_\mathrm{E}$ has 2 significant
  figures, so 2 or 3 are sensible.)}
\checkstep{$U$ is negative \cmark. Using the altitude as $r$ gives
  \qty{-3.39e11}{J}, about 8.5 times too big
  ($\num{7.221e6} \div \num{8.50e5} = 8.5$).}
\end{workedexample}

\begin{workedexample}{A research rocket: exact $\Delta U$ or $mg\Delta h$?}\label{ch5:we:rocket}
A research rocket carries a \qty{250}{kg} payload straight up. At the top, the
payload stops for a moment \qty{1500}{km} above Earth's surface. (a) Calculate
its gain in gravitational potential energy. (b) A student uses
$\Delta U = mg\Delta h$. Find the student's value and explain why it is too big.
\given{$m = \qty{250}{kg}$, \ $r_1 = r_\mathrm{E} = \qty{6.371e6}{m}$, \
  $r_2 = \qty{6.371e6}{m} + \qty{1.500e6}{m} = \qty{7.871e6}{m}$, \
  $\Delta h = \qty{1.500e6}{m}$, \ $g = \qty{9.8}{m.s^{-2}}$, \
  $G = \qty{6.67e-11}{N.m^2.kg^{-2}}$, \ $M_\mathrm{E} = \qty{6.0e24}{kg}$}
\find{(a) $\Delta U$ \quad (b) $mg\Delta h$, and the reason it is different}
\formula{$\Delta U = GMm\left(\dfrac{1}{r_1} - \dfrac{1}{r_2}\right)$ (not printed
  on the sheet: it is $U_2 - U_1$) \quad and \quad $\Delta U = mg\Delta h$}
\substitute{\eqlines{GMm &= (\num{6.67e-11})(\qty{6.0e24}{kg})(\qty{250}{kg}) = \qty{1.0005e17}{J.m}\\
  \frac{1}{r_1} - \frac{1}{r_2} &= \frac{1}{\qty{6.371e6}{m}} - \frac{1}{\qty{7.871e6}{m}}\\
     &= \qty{1.5696e-7}{m^{-1}} - \qty{1.2705e-7}{m^{-1}}\\
     &= \qty{2.9913e-8}{m^{-1}}\\
  \Delta U &= (\qty{1.0005e17}{J.m})(\qty{2.9913e-8}{m^{-1}}) = \qty{2.993e9}{J}\\
  mg\Delta h &= (\qty{250}{kg})(\qty{9.8}{m.s^{-2}})(\qty{1.500e6}{m}) = \qty{3.675e9}{J}}}
\answer{(a) \final{$\Delta U = \qty{2.99e9}{J}$} \quad (b) \final{$\qty{3.68e9}{J}$}.
  It is $\dfrac{3.675 - 2.993}{2.993} \times 100\,\% = 23\,\%$ too big.
  $mg\Delta h$ uses $g = \qty{9.8}{m.s^{-2}}$ all the way up, but $g$ falls to
  about \qty{6.5}{m.s^{-2}}. The real force is smaller, so less work is needed.}
\checkstep{$\Delta U > 0$ because the payload moved away from Earth \cmark.}
\end{workedexample}

\begin{mistakes}
\begin{itemize}
  \item \xmark\ Using the altitude as $r$. \cmark\ $r = r_\mathrm{E} + \text{altitude}$, in metres.
  \item \xmark\ Dropping the minus sign. \cmark\ $U$ is always negative.
  \item \xmark\ $\Delta U = \dfrac{GMm}{r_2 - r_1}$ or $mg\Delta h$ for orbits.
        \cmark\ Find $U_2 - U_1$.
\end{itemize}
\end{mistakes}

\begin{examtip}
``Explain why gravitational potential energy is negative'' (2--3 marks). Write
three linked points:
\begin{enumerate}
  \item We choose $U = 0$ at infinity, where gravity has died away to zero.
  \item Moving away needs work against gravity, so $U$ increases as $r$ increases.
  \item $U$ reaches 0 only at infinity. So it is negative at every real distance.
\end{enumerate}
In calculations, write $r = r_\mathrm{E} + \text{altitude}$ as its own line. It
often earns a mark.
\end{examtip}

\begin{quickcheck}
  \item A \qty{95}{kg} astronaut is in a space station \qty{410}{km} above Earth's
        surface. Calculate the astronaut's gravitational potential energy.
  \item Explain why the gravitational potential energy of a satellite is
        negative.\label{ch5:qc:negative}
  \item A \qty{500}{kg} satellite moves from $r = \qty{7.0e6}{m}$ to
        $r = \qty{9.0e6}{m}$. Calculate $\Delta U$ and the work done by gravity.
  \item A probe's distance from a planet's centre doubles. What happens to its
        $U$? Does $U$ increase or decrease?
\end{quickcheck}

%% ==================================================================
\section{Kinetic, potential and total energy in orbit}\label{ch5:sec:orbit}
\covers{M5.IQ3.05, M5.IQ3.04}

\begin{bigidea}
A satellite has kinetic energy because it moves, and potential energy because it
is in Earth's field. In a circular orbit the total is $U + K = -\dfrac{GMm}{2r}$.
A lower satellite moves faster, but it has \emph{less} total energy.
\end{bigidea}

\begin{keywords}
  \kw{Total mechanical energy ($E$)}{kinetic plus potential energy, $E = U + K$. Unit: J.}
  \kw{Bound}{trapped by gravity, so the object cannot get away. This happens when $E < 0$.}
  \kw{Perihelion, aphelion}{the closest and farthest points of an orbit around the Sun.}
\end{keywords}

\begin{refresher}{Orbital speed (\Cref{ch4:sec:orbits})}
In a circular orbit, gravity provides the centripetal force (the net force
towards the centre that keeps an object moving in a circle):
$\dfrac{GMm}{r^2} = \dfrac{mv^2}{r}$. Cancel $m$ and one $r$ from both sides:
$v^2 = \dfrac{GM}{r}$. Take the square root: $v = \sqrt{\dfrac{GM}{r}}$. This is
\emph{not} on the sheet. At \qty{400}{km} above Earth,
$v = \qty{7.69e3}{m.s^{-1}}$.
\end{refresher}

\subsection{Three energies of a satellite}

Remember the gravity well. A satellite in a higher orbit sits further up the
well. Getting there needs energy, like climbing stairs. Up there, gravity is
weaker, so a slower speed is enough to stay in orbit. So a higher orbit means a
slower satellite but more total energy.

Put $v^2 = \dfrac{GM}{r}$ into $K = \dfrac{1}{2}mv^2$:
$K = \dfrac{1}{2}m \times \dfrac{GM}{r} = \dfrac{GMm}{2r}$.

The \term{total mechanical energy} is the kinetic and potential energy added
together: $E = U + K$. For a circular orbit, $E = -\dfrac{GMm}{2r}$ (derived in
the card below).

``Mechanical'' means we count only $K$ and $U$. We do not count other energy,
such as the electrical energy stored in the satellite's batteries.

Compare the formulas. $K = \dfrac{GMm}{2r}$ and $E = -\dfrac{GMm}{2r}$, so
$E = -K$. Also $U = -\dfrac{GMm}{r} = 2 \times \left(-\dfrac{GMm}{2r}\right)$,
so $U = 2E$. These two shortcuts work for every circular orbit.

\Cref{ch5:fig:KUE} and \cref{ch5:tab:trends} show what happens in a bigger orbit.
The satellite slows down, but its total energy \emph{rises}. So a higher orbit
needs energy from the engines.

\begin{sidebyside}[0.53]
  \centering
  \begin{tikzpicture}
    \begin{axis}[hsc, width=0.94\linewidth, height=0.74\linewidth,
        xmin=1, xmax=6, ymin=-66, ymax=36, xtick={1,2,...,6},
        ytick={-60,-40,-20,0,20},
        xlabel={$r$ (in units of $r_\mathrm{E}$)},
        ylabel={energy of \qty{1}{kg} (\si{MJ})}]
      \draw[hscink!70, line width=0.8pt] (axis cs:1,0) -- (axis cs:6,0);
      \addplot[hscseriesA, domain=1:6, samples=80, line width=1.2pt, solid] {31.41/x};
      \addplot[hscseriesB, domain=1:6, samples=80, line width=1.2pt, dashed] {-62.82/x};
      \addplot[hscseriesC, domain=1:6, samples=80, line width=1.2pt, densely dotted] {-31.41/x};
      \node[font=\small, text=hscseriesA, anchor=south] at (axis cs:4.2,7.9) {$K$};
      \node[font=\small, text=hscseriesC, anchor=north] at (axis cs:2.5,-14.0) {$E$};
      \node[font=\small, text=hscseriesB, anchor=north] at (axis cs:2.5,-26.6) {$U$};
      \draw[guide] (axis cs:2,-31.41) -- (axis cs:2,15.70);
      \fill[hscseriesA] (axis cs:2,15.70) circle (1.6pt);
      \fill[hscseriesC] (axis cs:2,-15.70) circle (1.6pt);
      \fill[hscseriesB] (axis cs:2,-31.41) circle (1.6pt);
    \end{axis}
  \end{tikzpicture}
  \captionof{figure}{$K$, $U$ and $E$ for \qty{1}{kg} in a circular orbit around
    Earth. At $r = 2r_\mathrm{E}$ (dots): $K = +\qty{15.7}{MJ}$,
    $E = \qty{-15.7}{MJ}$, $U = \qty{-31.4}{MJ}$, so $E = -K$ and $U = 2E$.}\label{ch5:fig:KUE}
\nextside
  \captionof{table}{Moving to a bigger circular orbit.}\label{ch5:tab:trends}
  \centering\small
  \begin{tabular}{@{}L{31mm}L{33mm}@{}}
    \toprule
    \colhead{Quantity} & \colhead{When $r$ increases}\\
    \midrule
    speed $v$ & decreases\\
    force of gravity $F$ & decreases ($F \propto 1/r^2$)\\
    kinetic energy $K$ & decreases\\
    potential energy $U$ & increases (less negative)\\
    total energy $E$ & increases (less negative)\\
    period $T$ & increases\\
    \bottomrule
  \end{tabular}

  \raggedright\vspace{3mm}
  \textbf{Bound or not?} If $E < 0$, the object is \term{bound}: it orbits or
  falls back. If $E \ge 0$, it escapes. So a negative \emph{total} energy does
  mean something: the object is trapped.
\end{sidebyside}

\subsection{Energy in elliptical orbits (Kepler's laws)}
\covers{M5.IQ3.04, M5.IQ2.04}

\begin{refresher}{Kepler's three laws (\cref{ch4:sec:kepler})}
First law: each planet moves in an ellipse (a stretched circle). The Sun is at
one focus, a special point inside the ellipse that is not its centre.
Second law: a line from the Sun to the planet sweeps out equal areas in equal
times. Third law: $\dfrac{r^3}{T^2} = \dfrac{GM}{4\pi^2}$ (on the sheet).
\end{refresher}

In a circular orbit, gravity is always at \ang{90} to the velocity, so
$W = Fs\cos\ang{90} = 0$. Gravity does no work, and the speed stays the same.

In an ellipse, $r$ changes, so $U$ changes. Gravity is the only force, so
$E = K + U$ stays the same. When $U$ goes down, $K$ goes up
(\cref{ch5:fig:ellipse}).

\begin{sidebyside}[0.50]
  \centering
  \begin{tikzpicture}[x=0.8cm, y=0.8cm]
    \draw[hscmuted!70, line width=0.8pt] (0,0) ellipse (2.4 and 2.078);
    \draw[->, hscmuted, line width=0.6pt] (60:2.75 and 2.40) arc (60:110:2.75 and 2.40);
    \fill[hscacc] (1.2,0) circle (5pt);
    \node[diagramlabel, font=\footnotesize, below=5pt] at (1.2,0) {Sun};
    \draw[velocity] (2.4,0) -- (2.4,1.3) node[right, text=hscvel, font=\footnotesize] {$v$};
    \fill[object] (2.4,0) circle (3.5pt);
    \draw[velocity] (-2.4,0) -- (-2.4,-0.433) node[left, text=hscvel, font=\footnotesize] {$v$};
    \fill[object] (-2.4,0) circle (3.5pt);
    \coordinate (P) at (-1.2,-1.80);
    \draw[velocity] (P) -- ++(0.894,-0.447) node[below, text=hscvel, font=\footnotesize] {$v$};
    \draw[force] (P) -- ++(0.80,0.60) node[above left, text=hscforce, inner sep=1pt, font=\footnotesize] {$F$};
    \fill[object] (P) circle (3.5pt);
    \node[diagramlabel, font=\footnotesize, anchor=west] at (2.45,-0.35) {perihelion};
    \node[diagramlabel, font=\footnotesize, anchor=east] at (-2.45,0.35) {aphelion};
  \end{tikzpicture}
  \captionof{figure}{A planet in an elliptical orbit (motion anticlockwise). On
    the way in, part of $F$ points along $v$, so gravity speeds the planet
    up.}\label{ch5:fig:ellipse}
\nextside
  \begin{itemize}
    \item \textbf{Perihelion} (closest): $U$ is lowest, so $K$ is largest. The
          planet is fastest.
    \item \textbf{Aphelion} (farthest): $U$ is highest, so $K$ is smallest. The
          planet is slowest.
    \item Energy explains why the planet speeds up as it comes closer. This
          agrees with Kepler's second law. The line from the Sun to the planet
          sweeps out equal areas in equal times, so the planet moves fastest
          when it is closest.
  \end{itemize}
\end{sidebyside}

Kepler's third law gives $r$ when a question gives the period $T$ instead.

\begin{formulacard}[notonsheet]{Kinetic and total energy in a circular orbit}
  \cardformula{K = \frac{GMm}{2r} \qquad\qquad E = U + K = -\frac{GMm}{2r}}
  \begin{symbols}
    \sym{$K$}{kinetic energy of the satellite}{\si{J}}
    \sym{$E$}{total mechanical energy, $U + K$ ($U$: potential energy)}{\si{J}}
    \sym{$G$}{universal gravitational constant, \num{6.67e-11}}{\si{N.m^2.kg^{-2}}}
    \sym{$M$, $m$}{mass of the planet, mass of the satellite}{\si{kg}}
    \sym{$r$}{radius of the orbit, from the centre of the planet}{\si{m}}
  \end{symbols}
  \inwords{$K = \dfrac{GMm}{2r}$: a lower satellite moves faster, so it has more
    kinetic energy. $E = -\dfrac{GMm}{2r}$: a lower satellite has less (more
    negative) total energy; $E$ is always negative and is half of $U$.}
  \whentouse{\emph{circular} orbits only: comparing orbits and changing orbit.
    (The syllabus writes $U + K = -\dfrac{GMm}{2r}$.)}
  \whennottouse{an ellipse, a rocket, a falling object, or deciding whether an
    object is bound. Use $K = \tfrac{1}{2}mv^2$ and
    $E = \tfrac{1}{2}mv^2 - \dfrac{GMm}{r}$ at one point instead (next card).}
  \rearranged{$E = -K$ \qquad $U = 2E$ \qquad $r = \dfrac{GMm}{2K} = -\dfrac{GMm}{2E}$}
  \tensecond{\qty{1000}{kg} orbiting Earth at $r = \qty{7.0e6}{m}$:
    $K = \dfrac{(\num{6.67e-11})(\num{6.0e24})(1000)}{2 \times \num{7.0e6}} = \qty{2.86e10}{J}$,
    so $E = -K = \qty{-2.86e10}{J}$ and $U = 2E = \qty{-5.72e10}{J}$.}
  \derive{gravity is the centripetal force.
    \begin{align*}
      \frac{mv^2}{r} &= \frac{GMm}{r^2} && F_\mathrm{c} = F\\
      K = \tfrac{1}{2}mv^2 &= \frac{GMm}{2r} && \text{multiply both sides by $\tfrac{r}{2}$}\\
      E = U + K &= -\frac{2GMm}{2r} + \frac{GMm}{2r} = -\frac{GMm}{2r}
        && \text{add $U = -\tfrac{2GMm}{2r}$ (same denominator)}
    \end{align*}}
  \trap{circular orbits only. ``$E$ increases'' means it gets \emph{less} negative.}
\end{formulacard}

\begin{formulacard}[combine={$K = \frac{1}{2}mv^2$ and $U = -\frac{GMm}{r}$}]{Total energy at any point of a path}
  \cardformula{E = K + U = \frac{1}{2}mv^2 - \frac{GMm}{r}}
  \begin{symbols}
    \sym{$E$}{total mechanical energy}{\si{J}}
    \sym{$m$}{mass of the moving object}{\si{kg}}
    \sym{$v$}{speed of the object at that point}{\si{m.s^{-1}}}
    \sym{$G$}{universal gravitational constant, \num{6.67e-11}}{\si{N.m^2.kg^{-2}}}
    \sym{$M$}{mass of the planet or star}{\si{kg}}
    \sym{$r$}{distance from the centre of the planet or star at that point}{\si{m}}
  \end{symbols}
  \inwords{add the kinetic energy and the potential energy at one point of the path.}
  \whentouse{any path where gravity is the only force: an ellipse, a circle, or a
    projectile after launch. Also to decide whether an object is bound
    ($E < 0$) or escapes ($E \ge 0$).}
  \whennottouse{while an engine is firing or air drag is acting: then $E$ changes.}
  \rearranged{same $E$ at two points:
    $\tfrac{1}{2}v_1^2 - \dfrac{GM}{r_1} = \tfrac{1}{2}v_2^2 - \dfrac{GM}{r_2}$
    ($m$ cancels) \qquad $v_2^2 = v_1^2 - 2GM\left(\dfrac{1}{r_1} - \dfrac{1}{r_2}\right)$}
  \tensecond{\qty{1.0}{kg} at $r = \qty{1.0e7}{m}$ from Earth's centre, moving at
    \qty{7.0}{km.s^{-1}}: $K = \qty{2.45e7}{J}$ and $U = \qty{-4.00e7}{J}$, so
    $E = \qty{-1.55e7}{J}$. It is negative, so the object is bound.}
  \derive{add the two energies, both on the sheet.
    \begin{align*}
      E &= K + U && \text{total mechanical energy}\\
      &= \tfrac{1}{2}mv^2 + \left(-\frac{GMm}{r}\right) = \tfrac{1}{2}mv^2 - \frac{GMm}{r}
        && \text{put in $K$ and $U$}
    \end{align*}}
  \trap{use $v$ and $r$ at the \emph{same} point. Do not use $K = \dfrac{GMm}{2r}$
    here: that is for circles only.}
\end{formulacard}

\begin{workedexample}{Energy of an Earth-observation satellite}\label{ch5:we:observe}
A \qty{1600}{kg} Earth-observation satellite is in a circular orbit \qty{700}{km}
above Earth's surface. Calculate its kinetic, potential and total energy.
\given{$m = \qty{1600}{kg}$, \ $r = \qty{6.371e6}{m} + \qty{0.700e6}{m} = \qty{7.071e6}{m}$}
\find{$K$, $U$ and $E$}
\formula{$K = \dfrac{GMm}{2r}$, \quad $U = -\dfrac{GMm}{r}$, \quad $E = U + K$}
\substitute{\eqlines{GMm &= (\num{6.67e-11})(\qty{6.0e24}{kg})(\qty{1600}{kg}) = \qty{6.403e17}{J.m}\\
  K &= \frac{\qty{6.403e17}{J.m}}{2 \times \qty{7.071e6}{m}} = \qty{4.528e10}{J}\\
  U &= -\frac{\qty{6.403e17}{J.m}}{\qty{7.071e6}{m}} = \qty{-9.056e10}{J}\\
  E &= \qty{-9.056e10}{J} + \qty{4.528e10}{J} = \qty{-4.528e10}{J}}}
\answer{\final{$K = \qty{4.53e10}{J}$}, \final{$U = \qty{-9.06e10}{J}$},
  \final{$E = \qty{-4.53e10}{J}$}}
\checkstep{$E = -K$ and $U = 2E$ \cmark. Also $v = \sqrt{GM/r} = \qty{7.52e3}{m.s^{-1}}$
  gives $\tfrac{1}{2}mv^2 = \qty{4.53e10}{J}$, the same $K$.}
\end{workedexample}

\begin{toolkit}{Finding $r$ from the period: cube roots}
Multiply $\dfrac{r^3}{T^2} = \dfrac{GM}{4\pi^2}$ by $T^2$: $r^3 = \dfrac{GMT^2}{4\pi^2}$.
Then take the cube root (the $\sqrt[3]{\ }$ key, or power $\tfrac{1}{3}$). Example:
$\sqrt[3]{\num{8.0e21}} = \num{2.0e7}$. Change $T$ to \emph{seconds} first.
\end{toolkit}

\begin{workedexample}{A probe orbiting Mars}\label{ch5:we:mars}
An \qty{850}{kg} probe moves in a circular orbit around Mars with a period of
\qty{2.00}{h}. The mass of Mars is \qty{6.42e23}{kg}.
(a) Calculate the radius of the orbit. (b) Calculate the probe's total energy.
(c) State its kinetic energy.
\given{$m = \qty{850}{kg}$, \ $M = \qty{6.42e23}{kg}$, \
  $T = 2.00 \times \qty{3600}{s} = \qty{7200}{s}$}
\find{(a) $r$ \quad (b) $E$ \quad (c) $K$}
\formula{(a) $r^3 = \dfrac{GMT^2}{4\pi^2}$ \quad (b) $E = -\dfrac{GMm}{2r}$ \quad (c) $K = -E$}
\substitute{\eqlines{GM &= (\num{6.67e-11})(\qty{6.42e23}{kg}) = \qty{4.282e13}{m^3.s^{-2}}\\
  (\qty{7200}{s})^2 &= \qty{5.184e7}{s^2}, \qquad 4\pi^2 = 4 \times 3.1416^2 = 39.48\\
  r^3 &= \frac{(\qty{4.282e13}{m^3.s^{-2}})(\qty{5.184e7}{s^2})}{39.48}
     = \frac{\qty{2.220e21}{m^3}}{39.48}\\
     &= \qty{5.623e19}{m^3}\\
  r &= \sqrt[3]{\qty{5.623e19}{m^3}} = \qty{3.831e6}{m}\\
  E &= -\frac{(\qty{4.282e13}{m^3.s^{-2}})(\qty{850}{kg})}{2 \times \qty{3.831e6}{m}}
     = \qty{-4.750e9}{J}}}
\answer{(a) \final{$r = \qty{3.83e6}{m}$}\\
  (b) \final{$E = \qty{-4.75e9}{J}$}\\
  (c) \final{$K = -E = \qty{4.75e9}{J}$}}
\checkstep{Mars' radius is about \qty{3.39e6}{m}, so the probe is about
  \qty{440}{km} up: a sensible low orbit \cmark. Also $v = \dfrac{2\pi r}{T}
  = \qty{3.34e3}{m.s^{-1}}$ gives $\tfrac{1}{2}mv^2 = \qty{4.75e9}{J}$ \cmark.}
\end{workedexample}

\begin{workedexample}{Speed at the far end of an elliptical orbit}\label{ch5:we:ellipse}
A \qty{1200}{kg} satellite moves in an elliptical orbit around Earth. At
perigee (closest point) it is $r_\mathrm{p} = \qty{7.0e6}{m}$ from Earth's
centre and moves at \qty{8.50}{km.s^{-1}}. At apogee (farthest point) it is
$r_\mathrm{a} = \qty{1.20e7}{m}$ from the centre. (a) Calculate its total
energy. (b) Use energy to find its speed at apogee.
\given{$m = \qty{1200}{kg}$, \ $v_\mathrm{p} = \qty{8.50}{km.s^{-1}} = \qty{8.50e3}{m.s^{-1}}$, \
  $r_\mathrm{p} = \qty{7.0e6}{m}$, \ $r_\mathrm{a} = \qty{1.20e7}{m}$, \
  $GM = (\num{6.67e-11})(\qty{6.0e24}{kg}) = \qty{4.002e14}{m^3.s^{-2}}$}
\find{(a) $E$ \quad (b) $v_\mathrm{a}$}
\formula{$E = \tfrac{1}{2}mv^2 - \dfrac{GMm}{r}$. Gravity is the only force, so
  $E$ is the same at apogee: $K_\mathrm{a} = E - U_\mathrm{a}$.}
\substitute{\eqlines{K_\mathrm{p} &= \tfrac{1}{2}(\qty{1200}{kg})(\qty{8.50e3}{m.s^{-1}})^2 = \qty{4.335e10}{J}\\
  U_\mathrm{p} &= -\frac{(\qty{4.002e14}{m^3.s^{-2}})(\qty{1200}{kg})}{\qty{7.0e6}{m}} = \qty{-6.861e10}{J}\\
  E &= \qty{4.335e10}{J} + (\qty{-6.861e10}{J}) = \qty{-2.526e10}{J}\\
  U_\mathrm{a} &= -\frac{(\qty{4.002e14}{m^3.s^{-2}})(\qty{1200}{kg})}{\qty{1.20e7}{m}} = \qty{-4.002e10}{J}\\
  K_\mathrm{a} &= E - U_\mathrm{a} = \qty{-2.526e10}{J} - (\qty{-4.002e10}{J}) = \qty{1.476e10}{J}\\
  v_\mathrm{a} &= \sqrt{\frac{2K_\mathrm{a}}{m}} = \sqrt{\frac{2 \times \qty{1.476e10}{J}}{\qty{1200}{kg}}}
     = \qty{4.96e3}{m.s^{-1}}}}
\answer{(a) \final{$E = \qty{-2.53e10}{J}$}\\
  (b) \final{$v_\mathrm{a} = \qty{4.96e3}{m.s^{-1}}$}}
\checkstep{$E < 0$, so the satellite is bound \cmark. It is slower at apogee:
  further away, $U$ is higher, so $K$ is lower, as Kepler's second law
  says \cmark. $K$ falls by $\qty{4.335e10}{J} - \qty{1.476e10}{J} = \qty{2.859e10}{J}$.
  $U$ rises by $\qty{-4.002e10}{J} - (\qty{-6.861e10}{J}) = \qty{2.859e10}{J}$:
  the same amount, so $E$ stays the same \cmark.}
\end{workedexample}

\begin{mistakes}
\begin{itemize}
  \item \xmark\ Total energy $= -\dfrac{GMm}{r}$. \cmark\ That is $U$; $E = -\dfrac{GMm}{2r}$.
  \item \xmark\ ``A lower orbit is faster, so it has more total energy.''
        \cmark\ $K$ is bigger, but $E$ is lower (more negative).
  \item \xmark\ $K = \dfrac{GMm}{2r}$ for an ellipse. \cmark\ Use $\tfrac{1}{2}mv^2$.
  \item \xmark\ $T$ in hours in Kepler's third law. \cmark\ Use seconds.
\end{itemize}
\end{mistakes}

\begin{examtip}
``Derive $U + K = -\dfrac{GMm}{2r}$'' (3 marks). One mark for each step in the
card above:
\begin{enumerate}
  \item $F_\mathrm{c} = F$ gives $v^2 = \dfrac{GM}{r}$.
  \item So $K = \tfrac{1}{2}mv^2 = \dfrac{GMm}{2r}$.
  \item Add $U = -\dfrac{GMm}{r}$ to get $E = -\dfrac{GMm}{2r}$.
\end{enumerate}
Ratio multiple-choice questions (``$r$ doubles, so \ldots'') are common. Look at
where $r$ sits in each formula (Maths toolkit in \cref{ch5:sec:gpe}). Remember
that $U$ and $E$ are negative.
\end{examtip}

\begin{quickcheck}
  \item A \qty{300}{kg} satellite is in a circular orbit \qty{2000}{km} above
        Earth's surface. Calculate its $K$, $U$ and $E$.
  \item Earth (mass \qty{6.0e24}{kg}) moves around the Sun (mass
        \qty{2.0e30}{kg}) in a nearly circular orbit of radius \qty{1.5e11}{m}.
        Calculate Earth's gravitational potential energy and its total
        energy in this orbit.\label{ch5:qc:sun}
  \item A \qty{4.2e5}{kg} space station moves in a circular orbit around an
        unknown planet at \qty{6.8}{km.s^{-1}}. Neither the planet's mass nor
        the orbit's radius is known. Calculate the station's total
        energy.\label{ch5:qc:station}
  \item A satellite moves to a circular orbit with twice the radius. For each of
        $v$, $K$, $U$ and $E$, state whether it increases or decreases, and by
        what factor.
  \item Explain why gravity does no work on a satellite in a circular orbit, but
        does do work on a satellite in an elliptical orbit.
\end{quickcheck}

%% ==================================================================
\section{Escape velocity}\label{ch5:sec:escape}
\covers{M5.IQ3.05}

\begin{bigidea}
Throw a ball faster and it goes higher before it falls back. At or above one
special launch speed, the escape velocity, it never comes back. At exactly this
speed, its total energy is zero.
\end{bigidea}

\begin{keywords}
  \kw{Escape velocity ($v_\mathrm{esc}$)}{the smallest launch speed that lets a
    projectile get away from a planet for ever.}
  \kw{Projectile}{an object that gets all its speed at the start, then moves with
    only gravity acting (no engine).}
\end{keywords}

You met Newton's cannon in \cref{ch4:sec:orbits} (\cref{ch4:fig:cannon}): a
cannon on a very tall mountain, above the air, fires sideways faster and faster.
\Cref{ch5:fig:cannon} adds the fastest paths, which escape. The paths have their
true shapes. Only the mountain is drawn much too tall, so that you can see it.

\begin{sidebyside}[0.52]
  \centering
  \begin{tikzpicture}[scale=1.0]
    \fill[hscseriesA!12] (0,0) circle (1);
    \draw[hscseriesA!60, line width=0.8pt] (0,0) circle (1);
    \fill[hscmuted!45] (-0.20,0.98) -- (0,1.3) -- (0.20,0.98) -- cycle;
    \draw[hscmuted, line width=1.1pt, dashed, domain=0:44.62, samples=50, variable=\t]
      plot ({0.637/(1-0.51*cos(\t))*sin(\t)}, {0.637/(1-0.51*cos(\t))*cos(\t)});
    \draw[hscseriesA, line width=1.1pt] (0,0) circle (1.3);
    \draw[hscseriesC, line width=1.1pt, domain=0:360, samples=150, variable=\t]
      plot ({1.573/(1+0.21*cos(\t))*sin(\t)}, {1.573/(1+0.21*cos(\t))*cos(\t)});
    \draw[hscseriesB, line width=1.1pt, domain=0:108, samples=90, variable=\t]
      plot ({2.6/(1+cos(\t))*sin(\t)}, {2.6/(1+cos(\t))*cos(\t)});
    \draw[hscforce, line width=1.1pt, densely dashed, domain=0:96, samples=90, variable=\t]
      plot ({3.328/(1+1.56*cos(\t))*sin(\t)}, {3.328/(1+1.56*cos(\t))*cos(\t)});
    \fill[object] (0,1.3) circle (2.8pt);
    \node[diagramlabel, above, font=\footnotesize] at (0,1.36) {launch};
    \node[diagramlabel, font=\footnotesize] at (0,-0.25) {Earth};
    \draw[hscmuted, line width=0.4pt] (0.72,0.70) -- (0.55,0.30);
    \node[diagramlabel, font=\footnotesize, text=hscmuted] at (0.50,0.15) {1};
    \node[diagramlabel, font=\footnotesize, text=hscseriesA] at (-1.44,0) {2};
    \node[diagramlabel, font=\footnotesize, text=hscseriesC, anchor=north] at (0,-2.0) {3};
    \node[diagramlabel, font=\footnotesize, text=hscseriesB, anchor=north west] at (3.42,-1.05) {4};
    \node[diagramlabel, font=\footnotesize, text=hscforce, anchor=west] at (3.98,-0.42) {5};
  \end{tikzpicture}
  \captionof{figure}{Newton's cannon, firing sideways at five speeds (not to
    scale: the mountain is drawn far too tall).}\label{ch5:fig:cannon}
\nextside
  \begin{enumerate}
    \item slower than orbital speed: falls back
    \item $v = \sqrt{GM/r}$: circle
    \item between $\sqrt{GM/r}$ and $v_\mathrm{esc}$: ellipse
    \item $v = v_\mathrm{esc}$: just escapes
    \item $v > v_\mathrm{esc}$: escapes with speed to spare
  \end{enumerate}
  Paths 1--3 have $E < 0$ (bound). Path 4 has $E = 0$; path 5 has $E > 0$.
\end{sidebyside}

The total energy decides the path. To escape with the \emph{smallest} speed, the
object should reach infinity with no speed left. There $K = 0$ and $U = 0$, so
$E = 0$. Gravity is the only force, so $E$ stays the same all the way. So $E$
must also be zero at launch (path 4).

\begin{formulacard}[notonsheet]{Escape velocity}
  \cardformula{v_\mathrm{esc} = \sqrt{\frac{2GM}{r}}}
  \begin{symbols}
    \sym{$v_\mathrm{esc}$}{escape velocity. It is really a speed: the direction
      does not matter}{\si{m.s^{-1}}}
    \sym{$G$}{universal gravitational constant, \num{6.67e-11}}{\si{N.m^2.kg^{-2}}}
    \sym{$M$}{mass of the planet, moon or star}{\si{kg}}
    \sym{$r$}{distance from the centre where the object starts}{\si{m}}
  \end{symbols}
  \inwords{the smallest launch speed that lets a projectile get away for ever.}
  \whentouse{a projectile with no engine after launch, starting at distance $r$
    (on the surface or in orbit).}
  \whennottouse{a rocket whose engine keeps running, or when air resistance matters.}
  \rearranged{$M = \dfrac{v_\mathrm{esc}^2\, r}{2G}$ \qquad
    $r = \dfrac{2GM}{v_\mathrm{esc}^2}$ \qquad kinetic energy needed $= \dfrac{GMm}{r}$}
  \tensecond{Earth's surface:
    $v_\mathrm{esc} = \sqrt{\dfrac{2(\num{6.67e-11})(\num{6.0e24})}{\num{6.371e6}}}
    = \qty{1.12e4}{m.s^{-1}}$ (about \qty{40000}{km.h^{-1}}).}
  \derive{energy is conserved (gravity is the only force). In the smallest case,
    $K = 0$ and $U = 0$ at infinity, so $K + U = 0$ at the start.
    \begin{align*}
      \tfrac{1}{2}mv^2 - \frac{GMm}{r} &= 0 && \text{put in $K$ and $U$}\\
      \tfrac{1}{2}mv^2 &= \frac{GMm}{r} && \text{add $\tfrac{GMm}{r}$ to both sides}\\
      \tfrac{1}{2}v^2 &= \frac{GM}{r} && \text{divide both sides by $m$ ($m$ cancels)}\\
      v^2 &= \frac{2GM}{r} && \text{multiply both sides by 2}\\
      v_\mathrm{esc} &= \sqrt{\frac{2GM}{r}} && \text{take the square root}
    \end{align*}}
  \trap{do not lose the 2. Without it you get the orbital speed $\sqrt{GM/r}$.}
\end{formulacard}

Three facts follow from the derivation above:
\begin{itemize}
  \item $m$ cancels, so a pebble and a spacecraft need the same speed. (The
        heavier one needs more kinetic energy, $K = \dfrac{GMm}{r}$.) The launch
        direction does not matter either (energy has no direction), if the path
        misses the planet.
  \item At the same $r$,
        $v_\mathrm{esc} = \sqrt{\dfrac{2GM}{r}} = \sqrt{2}\times\sqrt{\dfrac{GM}{r}}
        = \sqrt{2}\times v_\text{orbital}$. A satellite in orbit already has
        $K = \dfrac{GMm}{2r}$, so it needs only $\dfrac{GMm}{2r}$ more kinetic
        energy.
  \item A rocket need not reach $v_\mathrm{esc}$: its engines keep adding energy,
        and higher up $v_\mathrm{esc}$ is smaller. It escapes once $K + U$ reaches zero.
\end{itemize}

\begin{workedexample}{Escaping from the Moon}\label{ch5:we:moon}
The Moon has a mass of \qty{7.35e22}{kg} and a radius of \qty{1.737e6}{m}.
Calculate the escape velocity from its surface and compare it with Earth's
(\qty{1.12e4}{m.s^{-1}}).
\given{$M = \qty{7.35e22}{kg}$, \ $r = \qty{1.737e6}{m}$}
\find{$v_\mathrm{esc}$}
\formula{$v_\mathrm{esc} = \sqrt{\dfrac{2GM}{r}}$}
\substitute{\eqlines{v_\mathrm{esc} &= \sqrt{\frac{2(\qty{6.67e-11}{N.m^2.kg^{-2}})(\qty{7.35e22}{kg})}{\qty{1.737e6}{m}}}\\
  &= \sqrt{\qty{5.645e6}{m^2.s^{-2}}} = \qty{2376}{m.s^{-1}}}}
\answer{\final{$v_\mathrm{esc} = \qty{2.38e3}{m.s^{-1}}$}. Earth's value is
  $\dfrac{\qty{1.12e4}{m.s^{-1}}}{\qty{2.38e3}{m.s^{-1}}} = 4.7$ times bigger, so
  the Moon's value is about $\tfrac{1}{5}$ of Earth's.}
\checkstep{The Moon's gravity is much weaker, so a smaller value makes sense. This
  is why the small Apollo ascent engine could lift the astronauts off the Moon
  into lunar orbit.}
\end{workedexample}

\begin{workedexample}{Leaving a parking orbit}\label{ch5:we:parking}
A \qty{1500}{kg} space probe is in a circular parking orbit (a low, waiting
orbit) \qty{600}{km} above Earth. Its engine fires briefly to send it away from
Earth for ever.
(a) Calculate the orbital speed and the escape velocity at this height.
(b) Calculate the smallest extra energy the engine must supply.
\given{$m = \qty{1500}{kg}$, \ $r = \qty{6.371e6}{m} + \qty{0.600e6}{m} = \qty{6.971e6}{m}$}
\find{(a) $v_\text{orbital}$ and $v_\mathrm{esc}$ \quad (b) the extra energy}
\formula{$v_\text{orbital} = \sqrt{\dfrac{GM}{r}}$, \ $v_\mathrm{esc} = \sqrt{\dfrac{2GM}{r}}$.
  To escape, $E$ must rise from its orbit value to 0, so extra energy
  $= E_\text{final} - E_\text{start} = 0 - \left(-\dfrac{GMm}{2r}\right) = \dfrac{GMm}{2r}$.}
\substitute{\eqlines{\frac{GM}{r} &= \frac{(\num{6.67e-11})(\qty{6.0e24}{kg})}{\qty{6.971e6}{m}}
     = \qty{5.741e7}{m^2.s^{-2}}\\
  v_\text{orbital} &= \sqrt{\qty{5.741e7}{m^2.s^{-2}}} = \qty{7577}{m.s^{-1}}\\
  v_\mathrm{esc} &= \sqrt{2 \times \qty{5.741e7}{m^2.s^{-2}}} = \qty{10715}{m.s^{-1}}\\
  \frac{GMm}{2r} &= \frac{GM}{r}\times\frac{m}{2}
     = (\qty{5.741e7}{m^2.s^{-2}})(\qty{750}{kg}) = \qty{4.306e10}{J}\\
     &\quad\text{(reuse $\tfrac{GM}{r}$ from the first line)}}}
\answer{(a) \final{$v_\text{orbital} = \qty{7.58e3}{m.s^{-1}}$}, \final{$v_\mathrm{esc} = \qty{1.07e4}{m.s^{-1}}$}\\
  (b) \final{$\qty{4.31e10}{J}$}}
\checkstep{Use the two speeds:
  $\tfrac{1}{2}m(v_\mathrm{esc}^2 - v_\text{orbital}^2)$
  $= \tfrac{1}{2}(1500)(\num{1.1482e8} - \num{5.741e7}) = \qty{4.31e10}{J}$
  \cmark. This is less than half of what a projectile needs
  from the ground ($GMm/r_\mathrm{E} = \qty{9.42e10}{J}$). The probe is already high up and already
  moving fast.}
\end{workedexample}

\begin{mistakes}
\begin{itemize}
  \item \xmark\ Forgetting the 2: $\sqrt{GM/r}$ is the orbital speed.
        \cmark\ $v_\mathrm{esc} = \sqrt{2GM/r}$.
  \item \xmark\ Deriving $v_\mathrm{esc}$ from $F_\mathrm{c} = F$. \cmark\ Use energy: $K + U = 0$.
  \item \xmark\ Using the surface value for a probe in orbit. \cmark\ Use the $r$ where it starts.
  \item \xmark\ ``A rocket must reach \qty{11.2}{km.s^{-1}}.'' \cmark\ Only a projectile must.
\end{itemize}
\end{mistakes}

\begin{examtip}
``Derive the escape velocity'' (2--3 marks):
\begin{enumerate}
  \item State that energy is conserved, because gravity is the only force.
  \item In the smallest case, $K = 0$ and $U = 0$ at infinity, so
        $\tfrac{1}{2}mv^2 - \dfrac{GMm}{r} = 0$ at launch.
  \item Show every step of the algebra.
\end{enumerate}
``Explain why a rocket can escape below $v_\mathrm{esc}$'' (3 marks):
\begin{enumerate}
  \item Escape velocity is for a projectile (no engine after launch).
  \item A rocket's engines keep adding energy as it climbs.
  \item $v_\mathrm{esc}$ is smaller at larger $r$.
  \item The rocket escapes when $K + U$ reaches zero.
\end{enumerate}
\end{examtip}

\begin{quickcheck}
  \item Mars has a mass of \qty{6.42e23}{kg} and a radius of \qty{3.39e6}{m}.
        Calculate the escape velocity from its surface.
  \item\label{ch5:qc:twoprobes} Two probes, of mass \qty{10}{kg} and
        \qty{1000}{kg}, are launched from Earth's surface. Do they need the same
        speed to escape? The same kinetic energy? Explain.%
  \item A rocket's engines fire for several minutes as it climbs. Its speed
        never goes above \qty{9}{km.s^{-1}}, but its payload still escapes from
        Earth. Use total energy to explain how.
  \item Calculate the smallest kinetic energy a \qty{40}{kg} probe needs to escape
        from the Moon's surface (data in \cref{ch5:we:moon}).
\end{quickcheck}

%% ==================================================================
\section{Moving between orbits}\label{ch5:sec:transfer}
\covers{M5.IQ3.05, M5.IQ3.03}

\begin{bigidea}
To move a satellite to a higher orbit, its engines must add energy $\Delta E$.
Surprisingly, it ends up moving more slowly: $K$ falls by $\Delta E$ and $U$ rises
by $2\Delta E$. Air drag does the opposite: it lowers the orbit.
\end{bigidea}

\begin{keywords}
  \kw{Hohmann transfer}{a fuel-saving way to change orbit, using two short engine
    firings (burns).}
  \kw{Transfer orbit}{the elliptical path between the two circular orbits.}
  \kw{Perigee, apogee}{the closest and farthest points of an orbit around Earth.}
  \kw{Parking orbit}{a low circular orbit where a spacecraft waits before its
    next engine burn.}
  \kw{Orbital decay}{a slow fall to lower orbits, caused by air drag.}
\end{keywords}

\subsection{The energy needed to change orbit}

Going to a higher orbit is like climbing higher up the gravity well, so it needs
energy. A bigger orbit has a less negative total energy, $E = -\dfrac{GMm}{2r}$.
So the engines must supply $\Delta E = E_2 - E_1$ (formula card below).

In every circular orbit $K = -E$, so
$\Delta K = K_2 - K_1 = -E_2 + E_1 = -\Delta E$. Also $U = 2E$, so
$\Delta U = 2\Delta E$. Check: $\Delta K + \Delta U = -\Delta E + 2\Delta E =
\Delta E$ \cmark.

\subsection{The Hohmann transfer}

The satellite follows the steps in \cref{ch5:fig:hohmann}. The speeds are for a
move from \qty{250}{km} up to geostationary orbit (GEO).
\begin{enumerate}
  \item \textbf{Burn 1} at perigee. A \term{burn} is a short time when the
        engine fires. Point the engine so that it pushes the satellite
        \emph{forwards}, in its direction of motion (the hot gas shoots out
        backwards). The speed rises from 7.77 to \qty{10.22}{km.s^{-1}}, and the
        circle stretches into an ellipse.
  \item \textbf{Coast} for half the ellipse, with no engine. Gravity slows the
        satellite as it moves away, turning $K$ into $U$; $E$ stays the same. At
        apogee it moves at only \qty{1.60}{km.s^{-1}}.
  \item \textbf{Burn 2} at apogee: push the satellite forwards again, up to the
        circular speed there, \qty{3.08}{km.s^{-1}}.
\end{enumerate}
Both burns speed the satellite up. But it ends up \emph{slower} than it started,
because it lost much more speed while coasting ``uphill''.

To move \emph{down} to a lower orbit, do the opposite. Push the satellite
backwards, against its motion, to remove energy. $E$ falls, the orbit gets
smaller, and the satellite ends up \emph{faster}. This is like orbital decay
(below), but controlled.

\begin{figure}
  \centering
  \begin{tikzpicture}[x=1.15cm, y=1.15cm, baseline=(current bounding box.center)]
    \fill[hscseriesA!15] (0,0) circle (0.55);
    \draw[hscseriesA!60, line width=0.8pt] (0,0) circle (0.55);
    \node[diagramlabel, font=\footnotesize] at (0,0) {Earth};
    \draw[hscmuted, line width=0.9pt] (0,0) circle (1.0);
    \draw[hscmuted, line width=0.9pt] (0,0) circle (2.3);
    \draw[hscseriesC, line width=1.2pt, domain=0:180, samples=80, variable=\t]
      plot ({-0.65+1.65*cos(\t)}, {1.5166*sin(\t)});
    \draw[hscseriesC!60, line width=0.7pt, dotted, domain=180:360, samples=60, variable=\t]
      plot ({-0.65+1.65*cos(\t)}, {1.5166*sin(\t)});
    \draw[hscseriesC, line width=1.2pt, ->, domain=40:72, samples=20, variable=\t]
      plot ({-0.65+1.65*cos(\t)}, {1.5166*sin(\t)});
    \draw[->, hscmuted, line width=0.7pt] (100:1.2) arc (100:140:1.2);
    \draw[guide] (0,0) -- (-60:1.0) node[pos=0.78, left, font=\footnotesize, text=hscmuted] {$r_1$};
    \draw[guide] (0,0) -- (-125:2.3) node[pos=0.55, left, font=\footnotesize, text=hscmuted] {$r_2$};
    \draw[accel] (1.0,0) -- (1.0,0.85) node[right, text=hscacc, font=\footnotesize, align=left] {burn 1};
    \draw[accel] (-2.3,0) -- (-2.3,-0.85) node[left, text=hscacc, font=\footnotesize, align=right] {burn 2};
    \fill[object] (1.0,0) circle (3pt);
    \fill[object] (-2.3,0) circle (3pt);
    \node[diagramlabel, font=\footnotesize, anchor=north west] at (1.0,-0.05) {perigee};
    \node[diagramlabel, font=\footnotesize, anchor=south east] at (-2.32,0.05) {apogee};
    \node[diagramlabel, font=\footnotesize, text=hscseriesC, align=center] at (-0.4,1.85) {transfer orbit};
  \end{tikzpicture}\hspace{5mm}%
  \begin{tikzpicture}[baseline=(current bounding box.center)]
    \begin{axis}[hsc, width=0.44\linewidth, height=0.40\linewidth,
        ybar, bar width=8pt, ymin=-66, ymax=36, grid=none,
        symbolic x coords={K, U, E}, xtick=data,
        xticklabels={$K$, $U$, $E$}, enlarge x limits=0.25,
        ytick={-60,-40,-20,0,20},
        ylabel={energy of \qty{1}{kg} (\si{MJ})},
        legend style={at={(0.5,1.03)}, anchor=south, legend columns=2,
          /tikz/every even column/.append style={column sep=4mm}},
        every axis plot/.append style={draw=none}]
      \addplot[fill=hscseriesA] coordinates {(K,30.22) (U,-60.44) (E,-30.22)};
      \addplot[fill=hscseriesB] coordinates {(K,4.73) (U,-9.46) (E,-4.73)};
      \legend{\qty{250}{km} orbit, GEO}
      \draw[hscink!70, line width=0.6pt] (rel axis cs:0,0.6471) -- (rel axis cs:1,0.6471);
    \end{axis}
  \end{tikzpicture}
  \caption{Left: a Hohmann transfer from a low orbit (radius $r_1$) to GEO
    (radius $r_2$). The small arrows show the motion (anticlockwise). Both burns
    push the satellite forwards, along its motion (not to scale: the real $r_2$
    is 6.4 times $r_1$). Dotted: the half of the ellipse the satellite does not
    use. Right: in GEO, $K$ is smaller, but $U$ and $E$ are higher (less
    negative).}\label{ch5:fig:hohmann}
\end{figure}

\subsection{From the ground, and satellites in practice}

On the launch pad, $K \approx 0$, so $E = U = -\dfrac{GMm}{r_\mathrm{E}}$.
(Earth's spin gives a rocket a small free speed; we leave it out.) The smallest
energy to reach a circular orbit is
\[ \Delta E = E_\text{orbit} - E_\text{start}
   = -\frac{GMm}{2r} - \left(-\frac{GMm}{r_\mathrm{E}}\right). \]
Example for \qty{1}{kg} going to \qty{250}{km}: $E_\text{start} = \qty{-62.8}{MJ}$
and $E_\text{orbit} = \qty{-30.2}{MJ}$, so
$\Delta E = -30.2 - (-62.8) = \qty{32.6}{MJ}$. To reach GEO it is
\qty{58.1}{MJ}, and to escape it is \qty{62.8}{MJ}. (Real rockets use far more:
they lift their fuel and push through air.)

So it costs much less energy to reach low Earth orbit (LEO) than geostationary
orbit (GEO). LEO is used by the International Space Station (ISS), Earth-imaging
satellites and Starlink. A GEO satellite is above the equator and goes round
once a day. TV satellites and weather satellites such as Himawari use GEO.
\Cref{ch4:sec:leogeo} compares the two orbits.

\textbf{Energy changes during a launch.} A common exam question asks you to
describe these in three stages.
\begin{enumerate}
  \item \textbf{On the pad:} energy is stored as chemical energy in the fuel.
        $K = 0$ and $U = -\dfrac{GMm}{r_\mathrm{E}}$.
  \item \textbf{While the engines burn:} the fuel's chemical energy becomes the
        rocket's extra $K$ and $U$. Some also becomes heat, sound and the kinetic
        energy of the exhaust gases.
  \item \textbf{After the engines stop:} the rocket coasts upwards and slows
        down. $K$ turns into $U$, and $K + U$ stays constant (ignoring air).
\end{enumerate}
It is $K + U$ that stays constant while coasting, not ``$K + U = 0$''.

\subsection{Orbital decay}

A LEO satellite still moves through a very thin layer of air. \term{Atmospheric
drag} is the push of this air against the satellite's motion. It turns some
mechanical energy into heat.

\begin{sidebyside}[0.36]
  \centering
  \begin{tikzpicture}[x=0.95cm, y=0.95cm]
    \fill[hscseriesA!15] (0,0) circle (0.9);
    \draw[hscseriesA!60, line width=0.8pt] (0,0) circle (0.9);
    \node[diagramlabel, font=\footnotesize] at (0,0) {Earth};
    \draw[hscmuted, line width=0.9pt, domain=0:1080, samples=240, variable=\t, ->]
      plot ({(2.0-0.6*\t/1080)*cos(\t)}, {(2.0-0.6*\t/1080)*sin(\t)});
    \coordinate (S) at ({1.975*cos(45)}, {1.975*sin(45)});
    \draw[velocity] (S) -- ++(135:0.95) node[above, text=hscvel, font=\footnotesize] {$v$};
    \draw[force] (S) -- ++(-45:0.55) node[right, text=hscforce, font=\footnotesize] {drag};
    \draw[force] (S) -- ($(S)!0.75cm!(0,0)$) node[anchor=north east, text=hscforce,
      font=\footnotesize, inner sep=1pt, fill=white, rounded corners=1pt] {gravity};
    \fill[object] (S) circle (3.5pt);
  \end{tikzpicture}
  \captionof{figure}{Drag makes the satellite spiral slowly inwards (spacing
    exaggerated).}\label{ch5:fig:decay}
\nextside
  Follow the energy:
  \begin{enumerate}
    \item Drag removes energy, so $E$ falls (more negative).
    \item $E = -\dfrac{GMm}{2r}$, so $r$ falls. The orbit gets lower.
    \item $E = -K$, so $K$ \emph{rises}: the satellite speeds up.
    \item $U = 2E$, so $U$ falls by twice the energy lost. Half of this fall in
          $U$ becomes heat. The other half becomes extra $K$.
  \end{enumerate}
  Lower orbits meet denser air, so decay gets faster, and the satellite finally
  burns up. GEO satellites have almost no air around them, so they hardly decay.
\end{sidebyside}

Drag slows the satellite for a moment, so it drops a little. As it drops,
gravity speeds it up by more than drag slowed it. So decay is caused by drag, not
by ``stronger gravity'': gravity on its own keeps an orbit steady for ever.

\begin{formulacard}[notonsheet]{Energy to change circular orbit}
  \cardformula{\Delta E = \frac{GMm}{2}\left(\frac{1}{r_1} - \frac{1}{r_2}\right)}
  \begin{symbols}
    \sym{$\Delta E$}{change in total energy (smallest work done by the engines)}{\si{J}}
    \sym{$r_1$, $r_2$}{radius of the starting orbit, radius of the new orbit}{\si{m}}
    \sym{$G$}{universal gravitational constant, \num{6.67e-11}}{\si{N.m^2.kg^{-2}}}
    \sym{$M$}{mass of the planet}{\si{kg}}
    \sym{$m$}{mass of the satellite}{\si{kg}}
  \end{symbols}
  \inwords{moving out to a bigger orbit needs energy; moving in gives energy out.}
  \whentouse{a satellite moves from one circular orbit to another (or decays).}
  \whennottouse{a launch from the ground: the start is not an orbit (see above).}
  \rearranged{$\Delta K = -\Delta E$ \qquad $\Delta U = 2\Delta E$}
  \tensecond{\qty{100}{kg} from $r_1 = \qty{7.0e6}{m}$ to $r_2 = \qty{1.4e7}{m}$ (Earth):
    $\dfrac{GMm}{2} = \qty{2.0e16}{J.m}$ and
    $\dfrac{1}{r_1} - \dfrac{1}{r_2} = \qty{7.14e-8}{m^{-1}}$, so
    $\Delta E = \qty{1.4e9}{J}$.}
  \derive{final minus initial total energy.
    \begin{align*}
      \Delta E &= E_2 - E_1 = \left(-\frac{GMm}{2r_2}\right) - \left(-\frac{GMm}{2r_1}\right)
        && \text{put in } E = -\tfrac{GMm}{2r}\\
      &= \frac{GMm}{2r_1} - \frac{GMm}{2r_2} && \text{minus a minus is a plus}\\
      &= \frac{GMm}{2}\left(\frac{1}{r_1} - \frac{1}{r_2}\right)
        && \text{take out the common factor } \tfrac{GMm}{2}
    \end{align*}}
  \trap{$\Delta E$ is not $\Delta U$. The engines supply only half of the rise in
    $U$; the other half comes from the fall in $K$.}
\end{formulacard}

\begin{workedexample}{From a parking orbit to geostationary orbit}\label{ch5:we:geo}
A \qty{2800}{kg} communications satellite is in a circular parking orbit
\qty{250}{km} above Earth. A Hohmann transfer moves it to geostationary orbit,
where the period is \qty{24.0}{h}.
(a) Show that the radius of geostationary orbit is about \qty{4.23e7}{m}.
(b) Calculate the smallest energy needed for the transfer.
(c) Calculate the change in kinetic energy, and explain why it is negative even
though both burns speed the satellite up.
\given{$m = \qty{2800}{kg}$, \ $r_1 = \qty{6.371e6}{m} + \qty{0.250e6}{m} = \qty{6.621e6}{m}$, \
  $T = 24.0 \times \qty{3600}{s} = \qty{86400}{s}$}
\find{(a) $r_2$ \quad (b) $\Delta E$ \quad (c) $\Delta K$}
\formula{(a) $r^3 = \dfrac{GMT^2}{4\pi^2}$ \quad
  (b) $\Delta E = \dfrac{GMm}{2}\left(\dfrac{1}{r_1} - \dfrac{1}{r_2}\right)$ \quad
  (c) $\Delta K = -\Delta E$}
\substitute{\eqlines{r_2^3 &= \frac{(\num{6.67e-11})(\qty{6.0e24}{kg})(\qty{86400}{s})^2}{4\pi^2}
     = \qty{7.567e22}{m^3}\\
  r_2 &= \sqrt[3]{\qty{7.567e22}{m^3}} = \qty{4.230e7}{m}\\
  \frac{GMm}{2} &= \frac{(\num{6.67e-11})(\num{6.0e24})(2800)}{2} = \qty{5.603e17}{J.m}\\
  \frac{1}{r_1} - \frac{1}{r_2} &= \frac{1}{\qty{6.621e6}{m}} - \frac{1}{\qty{4.230e7}{m}}
     = \qty{1.2739e-7}{m^{-1}}\\
  \Delta E &= (\qty{5.603e17}{J.m})(\qty{1.2739e-7}{m^{-1}}) = \qty{7.138e10}{J}}}
\answer{(a) \final{$r_2 = \qty{4.23e7}{m}$}\\
  (b) \final{$\Delta E = \qty{7.14e10}{J}$}\\
  (c) \final{$\Delta K = -\Delta E = \qty{-7.14e10}{J}$}. Between the burns the
  satellite coasts away from Earth, and gravity turns $K$ into $U$. The two
  burns add $\num{6.17e10} + \qty{0.97e10}{J} = \qty{7.14e10}{J}$ of $K$, but
  coasting removes \qty{1.43e11}{J}.}
\checkstep{$\Delta U = 2\Delta E = \qty{1.428e11}{J}$, and
  $\Delta K + \Delta U = \qty{-7.138e10}{J} + \qty{1.428e11}{J} = \qty{7.14e10}{J} = \Delta E$
  \cmark. Final speed \qty{3.08}{km.s^{-1}} is below the starting
  \qty{7.77}{km.s^{-1}} \cmark.}
\end{workedexample}

\begin{workedexample}{Orbital decay of a space telescope}\label{ch5:we:decay}
An \qty{11000}{kg} space telescope is in a circular orbit \qty{560}{km} above
Earth. Over several years, air drag lowers it to a circular orbit \qty{500}{km}
above Earth. (a) Calculate the mechanical energy changed into heat.
(b) Calculate the change in its speed.
\given{$m = \qty{11000}{kg}$, \
  $r_1 = \qty{6.371e6}{m} + \qty{0.560e6}{m} = \qty{6.931e6}{m}$, \
  $r_2 = \qty{6.371e6}{m} + \qty{0.500e6}{m} = \qty{6.871e6}{m}$}
\find{(a) the energy lost, $-\Delta E$ \quad (b) $\Delta v = v_2 - v_1$}
\formula{$\Delta E = \dfrac{GMm}{2}\left(\dfrac{1}{r_1} - \dfrac{1}{r_2}\right)$ \ and \
  $v = \sqrt{\dfrac{GM}{r}}$. Write $\tfrac{1}{r_1} - \tfrac{1}{r_2}$ as one
  fraction, $\tfrac{r_2 - r_1}{r_1 r_2}$ (Maths toolkit in \cref{ch5:sec:gpe}).
  The two radii are almost equal, so subtracting $\tfrac{1}{r_1}$ and
  $\tfrac{1}{r_2}$ on a calculator can lose digits.}
\substitute{\eqlines{\frac{GMm}{2} &= \frac{(\num{6.67e-11})(\num{6.0e24})(11000)}{2} = \qty{2.201e18}{J.m}\\
  \frac{1}{r_1} - \frac{1}{r_2} &= \frac{r_2 - r_1}{r_1 r_2}
     = \frac{\qty{-6.0e4}{m}}{(\qty{6.931e6}{m})(\qty{6.871e6}{m})}\\
     &= \qty{-1.2599e-9}{m^{-1}}\\
  \Delta E &= (\qty{2.201e18}{J.m})(\qty{-1.2599e-9}{m^{-1}}) = \qty{-2.773e9}{J}\\
  GM &= (\num{6.67e-11})(\qty{6.0e24}{kg}) = \qty{4.002e14}{m^3.s^{-2}}\\
  v_1 &= \sqrt{\frac{\qty{4.002e14}{m^3.s^{-2}}}{\qty{6.931e6}{m}}} = \qty{7598.7}{m.s^{-1}}\\
  v_2 &= \sqrt{\frac{\qty{4.002e14}{m^3.s^{-2}}}{\qty{6.871e6}{m}}} = \qty{7631.8}{m.s^{-1}}\\
  \Delta v &= v_2 - v_1 = \qty{7631.8}{m.s^{-1}} - \qty{7598.7}{m.s^{-1}} = +\qty{33.1}{m.s^{-1}}}}
\answer{(a) \final{$\qty{2.77e9}{J}$} becomes heat\\
  (b) \final{$\Delta v = +\qty{33}{m.s^{-1}}$}: the telescope speeds up.}
\checkstep{$\Delta K = -\Delta E = +\qty{2.77e9}{J}$, and
  $\tfrac{1}{2}m(v_2^2 - v_1^2) = \qty{2.77e9}{J}$ \cmark: lower and faster.}
\end{workedexample}

\begin{mistakes}
\begin{itemize}
  \item \xmark\ Using $\Delta U$ as the energy to change orbit. \cmark\ Use $\Delta E$.
  \item \xmark\ ``Slow the satellite down to make it go higher.'' \cmark\ Both
        burns push it forwards and speed it up.
  \item \xmark\ ``A higher satellite is faster.'' \cmark\ It is slower: $v = \sqrt{GM/r}$.
  \item \xmark\ ``Decay happens because gravity is stronger low down.'' \cmark\ Drag causes it.
\end{itemize}
\end{mistakes}

\begin{examtip}
``Explain how a satellite is moved from a low orbit to geostationary orbit''
(3--4 marks). Name the Hohmann transfer orbit. Then:
\begin{enumerate}
  \item A burn at perigee pushes the satellite forwards. Its orbit becomes an ellipse.
  \item While it coasts out, $K$ turns into $U$.
  \item A second forward burn at apogee gives the new circle.
\end{enumerate}
For full marks add the energy: $E$ and $U$ rise, $K$ falls. In ``assess''
questions on decay, give a clear judgement first (correct or incorrect). Then
give the physics: drag, energy loss, lower $r$, higher $v$.
\end{examtip}

\begin{quickcheck}
  \item A \qty{500}{kg} satellite moves from a circular orbit of radius
        \qty{8.0e6}{m} to one of radius \qty{2.4e7}{m}. Calculate $\Delta E$,
        $\Delta K$ and $\Delta U$.
  \item\label{ch5:qc:launch} Calculate the smallest energy to put a
        \qty{650}{kg} satellite from rest on Earth's surface into a circular
        orbit \qty{500}{km} high (ignore Earth's spin and air).
  \item Describe the energy changes of a rocket from the launch pad until it
        coasts upwards after its engines stop.\label{ch5:qc:launchenergy}
  \item Two engine burns push a satellite forwards, along its motion, into a
        higher circular orbit. Explain why it ends up moving more slowly.
  \item A student says: ``Air drag slows a satellite down. So a satellite whose
        orbit is decaying ends up moving more slowly.'' Assess this
        claim.\label{ch5:qc:assess}
\end{quickcheck}

%% ==================================================================
\section{Key experiment: launch speed, energy and escape}\label{ch5:sec:experiment}
\covers{M5.IQ3.05, M5.IQ3.04}

There is no hands-on practical here, so we investigate with a computer
simulation. The results below come from a simple Python model.

\begin{experiment}{Launch speed, total energy and escape (simulation)}
\aim{To find how the launch speed decides whether a probe falls back, orbits or
  escapes, and to test $v_\mathrm{esc} = \sqrt{2GM/r}$ using total energy.}
% layout QA: split 0.44 (was 0.50) and a wider 'path' column in the results table save two
% lines, so the chapter summary fits on the last page of the experiment
\apparatus{\begin{sidebyside}[0.44]
    \centering
    \begin{tikzpicture}[x=0.7cm, y=0.7cm]
      \draw[rounded corners=2mm, hscmuted, line width=0.8pt, fill=hscmuted!5] (-3.4,-1.45) rectangle (5.0,1.45);
      \fill[hscseriesA!15] (-1.6,0) circle (0.6);
      \draw[hscseriesA!60, line width=0.8pt] (-1.6,0) circle (0.6);
      \node[diagramlabel, font=\footnotesize] at (-1.6,0) {Earth};
      \draw[guide] (-1.6,0) -- node[below, font=\footnotesize, pos=0.75] {$r_0$} (-0.3,0);
      \draw[hscmuted, dashed, line width=0.6pt] (-1.6,0) circle (1.3);
      \fill[object] (-0.3,0) circle (3.2pt);
      \draw[velocity] (-0.3,0) -- (-0.3,1.05) node[right, text=hscvel, font=\footnotesize] {$v_0$};
      \draw[hscmuted, line width=0.6pt] (2.5,-1.45) -- (2.5,1.45);
      \node[font=\footnotesize, anchor=west, align=left, text=hscink] at (2.55,0.85) {launch\\[-2pt]speed};
      \draw[hscmuted, line width=1.2pt] (2.75,-0.05) -- (4.7,-0.05);
      \fill[hscvel] (3.7,-0.05) circle (2.5pt);
      \node[font=\footnotesize, anchor=west, align=left, text=hscink] at (2.55,-0.85) {shows $r$, $v$};
    \end{tikzpicture}
    \captionof{figure}{Simulation set-up: launch at \ang{90} to the radius.}\label{ch5:fig:sim}
  \nextside
    A computer orbit model, such as a spreadsheet or Python model. You type in
    the planet's mass, the start distance and the launch speed. (PhET
    ``Gravity and Orbits'' shows the path shapes, but it cannot set speeds this
    precisely.) Set the planet's mass to Earth's, \qty{6.0e24}{kg}
    (\cref{ch5:fig:sim}).
  \end{sidebyside}}
\method{\begin{enumerate}[beginpenalty=10000]
  \item Launch the probe sideways from $r_0 = \qty{1.00e7}{m}$ with speed $v_0$.
        Record the path and its closest and farthest distances.
  \item Repeat for $v_0$ from 5 to \qty{9}{km.s^{-1}} (small steps near 9). For each
        run calculate $E/m = \tfrac{1}{2}v_0^2 - GM/r_0$.
  \item For one ellipse, read $r$ and $v$ at the closest and farthest points.
        Calculate $K$, $U$ and $E$ for \qty{1}{kg} at each.
\end{enumerate}}
\results{\begin{center}\small
  \begin{tabular}{@{}ccL{48mm}cc@{}}
    \toprule
    $v_0$ (\si{km.s^{-1}}) & $E/m$ (\si{MJ.kg^{-1}}) & path & closest $r$ (\si{m}) & farthest $r$ (\si{m})\\
    \midrule
    5.00 & $-27.5$ & falls back, hits Earth & -- & \num{1.00e7}\\
    6.33 & $-20.0$ & circle (very nearly) & \num{1.00e7} & \num{1.00e7}\\
    7.00 & $-15.5$ & ellipse & \num{1.00e7} & \num{1.58e7}\\
    8.94 & $-0.06$ & very long ellipse (in the model, back after 2 years) & \num{1.00e7} & \num{6.9e9}\\
    8.95 & $+0.03$ & escapes & \num{1.00e7} & none\\
    \bottomrule
  \end{tabular}
  \end{center}
  In reality, beyond about \qty{1.5e9}{m} from Earth, the Sun's gravity takes
  control of the probe's path, so the \qty{8.94}{km.s^{-1}} probe would not come
  back. The model has only two masses.

  The \qty{7.00}{km.s^{-1}} ellipse, for \qty{1}{kg}: closest point, $v = \qty{7.00}{km.s^{-1}}$,
  $K = \qty{24.50}{MJ}$, $U = \qty{-40.02}{MJ}$, $E = \qty{-15.52}{MJ}$;
  farthest point, $v = \qty{4.43}{km.s^{-1}}$, $K = \qty{9.83}{MJ}$,
  $U = \qty{-25.35}{MJ}$, $E = \qty{-15.52}{MJ}$.}
\conclusion{At $v_0 = \sqrt{GM/r_0} = \qty{6.33}{km.s^{-1}}$ the path is a
  circle. The probe escapes only when $E \ge 0$. The smallest escape speed is
  between 8.94 and \qty{8.95}{km.s^{-1}}; theory gives
  $\sqrt{2GM/r_0} = \qty{8.95}{km.s^{-1}}$ (within 0.1\,\%). Around the ellipse,
  $E$ stays the same while $K$ and $U$ swap. The probe is fastest when closest,
  as Kepler's second law says.}
\significance{Total energy, not speed alone, decides whether an object is bound.
  Mission planners use this rule to choose launch speeds.}
\validity{\begin{itemize}[beginpenalty=10000]
  \item \textbf{Validity:} change only $v_0$; keep $r_0$, $M$ and the launch
        direction fixed. The model treats Earth and the probe as point masses
        (all their mass at one point) and ignores air. Far from Earth (beyond
        about \qty{1.5e9}{m}) the Sun's gravity matters, but the model leaves it out.
  \item \textbf{Reliability:} an identical run always gives an identical result,
        so repeating it tells you nothing new. Instead, repeat each run with half
        the \term{time step}: how far forward in time the computer jumps in each
        calculation. If the results stay the same, they are reliable.
  \item \textbf{Accuracy:} compare the escape speed found with $\sqrt{2GM/r_0}$
        (percentage difference). Too large a time step makes ``circles'' spiral.
\end{itemize}}
\end{experiment}

%% ==================================================================
% Chapter summary (added by the whole-book consistency edit: "How to use this guide" says
% that every chapter ends with a short summary). Every statement is taught above.
\begin{summarybox}[unbreakable]{\Cref{ch5:chap} in five lines}
\begin{itemize}
  \item $U = -\frac{GMm}{r}$ is zero at infinity and negative everywhere else. Use
        $\Delta U = mg\Delta h$ only near the surface.
  \item Circular orbit: $K = \frac{GMm}{2r}$ and $E = U + K = -\frac{GMm}{2r}$, so $E = -K$ and
        $U = 2E$.
  \item When gravity is the only force, $E = K + U$ stays constant. If $E < 0$ the object is
        bound; if $E \ge 0$ it escapes.
  \item Escape velocity: $K + U = 0$ gives $v_\mathrm{esc} = \sqrt{\frac{2GM}{r}}$
        (\qty{11.2}{km.s^{-1}} from the Earth's surface), whatever the mass.
  \item A higher orbit needs energy,
        $\Delta E = \frac{GMm}{2}\left(\frac{1}{r_1} - \frac{1}{r_2}\right)$: $E$ and $U$ rise,
        but $K$ falls. Drag lowers an orbit and speeds the satellite up.
\end{itemize}
\end{summarybox}
